\addcontentsline{toc}{chapter}{Introduction}
\chapter*{Introduction\label{chap_I}}
\epigraph{One should not pursue goals that are easily achieved. One must develop an instinct for what one can barely achieve through one's greatest efforts.}{---Albert Einstein}
\epigraph{Nothing will ever be attempted, if all possible objections must be first overcome.}{---`the artist', in Samuel Johnson's \textit{Rasselas}}
%This introduction is intended to help with the readability of the thesis, because its scope is very large, and its justification eventually does need to be considered and understood as a whole; and, I fully realise, because no one is currently of the same mindset as I am---which I've come to through mostly independent research over the past few years, in which time I have made much effort to always consider objectively whether any commonly-held beliefs as to the meaning of physical theory are truly justifiable, and have thus come to some fairly different opinions when analytical results seemed to suggest a more natural interpretation than one that conventional wisdom would suppose; so, without this heuristic outline of the thesis, which I will attempt to provide by synoptically covering not only what I think are its main results, but also what were some of the principal goals and developments that led to the realisation of those results along with the different point-of-view from which the whole thing works, it seems that it would be more difficult to follow the more detailed theoretical development, as the greater scope to which all of the points that are considered \textit{do} pertain might become less clear in the course of more rigorous investigation, or even discounted as untenable due to a bias which justification will eventually be argued against.

This thesis began from an idea, that there might be some way to describe not only the current accelerated rate of cosmic expansion, but actually the expansion of the Universe at all times, as a consequence of the basic metrical properties of physical reality. Specifically, the basic idea I had, which actually isn't what the eventual theory comes to describe, was that if our Universe had some essential characteristic which would cause it to naturally expand right from the Big Bang, which would be described by a pure cosmological constant, then at early times, when all the mass in the universe was close together, the mutual attraction would have a slowing effect on the rate of that intrinsic cosmic expansion, which would weaken as the universe did expand, until an inflection point was reached and $\Lambda$ began to dominate.

The problem with this idea is that it can't exactly be described by Friedman's equations, which are totally dominated by material densities near the big bang, so that no matter how colse the expansion of a Friedman-Lema\^{i}tre-Robertson-Walker (FLRW) universe might actually correspond to this description---e.g., as in the $\Lambda$CDM models---the description given by the standard cosmological model and the one I was thinking of could never be alike close to the Big Bang.\footnote{Note, that throughout the thesis, I use capitalisation to distinguish the names of (analytically uncertain) things that actually are, such as the Universe, Truth, and Physical Reality, from the things of physical theory, such as a universe, truth, or physical reality; e.g., in the end, after working out the physical properties of a certain universe, I propose that that is what our Universe actually is, with its big bang describing the Big Bang we believe to have taken place at the beginning of Time---which is something that also wants a more concrete analytical description.} This is why the classical FLRW `big bang' models \textit{require} that singular push to start things off, which might be subsequently helped along, e.g.~by an inflationary epoch, which would also smooth things out. 

Basically, this dynamical theory (and my idea wasn't much different, except that I preferred to think of $\Lambda$, and its role for expansion, as a basic metrical property) requires a dynamical origin for expansion, and I thought, instead, that there should be a way of describing it purely as a metrical property;\footnote{I should like to be absolutely clear, that I've never subscribed to the opinion that formal equivalence of the physical descriptions of two essentially different things means as much as if they Truly were identical, and have always felt that there is epistemological merit in trying to understand what things Really are. The contrary position is, as I understand it, one that was maintained by some of Johannes Kepler and Galileo Galilei's strongest adversaries, who eventually conceded the superiority of Kepler's heliocentric system only because it provided a more accurate description of the phenomena. But Kepler and Galileo did not have the luxury of those results when they fought their battles: they led the Copernican revolution in their day because a Truly heliocentric system would give rise more naturally to the phenomena. Of course, the fact that the empirically more accurate Keplerian theory played an essential role in Isaac Newton's eventual discovery of the far deeper theory of Universal gravitation, which allowed him to further refine the description of planetary orbits, is a clear demonstration of the fact that it really is worthwhile to search for essential Truth.} so, because I still thought there could be something more to this idea of expansion, I began looking into other options. This seemed reasonable enough to do, being that the empirical evidence does generally favour a pure cosmological constant, which, if it were a fundamental constant of Nature, should reasonably be thought to have an essential connection to the cosmic expansion that appears to be a fundamental property of the Universe.

Now, in addition to the fact that the standard cosmological model is incapable of providing an essential explanation of cosmic expansion---I mean, of the most basic cause of cosmic expansion, which it only describes as the consequence of an initial singularity---there are many other well-known problems with it, as the parameters that have been constrained through empirical modelling are, according to the most objective insights that can be made consistently within the dynamical theory, \textit{completely} unrealistic. To me, all of its problems suggested that there really could be something essentially wrong with the standard model, so that rather than lumping conjectures onto that model in order to smooth out the issues it is known to have, the essential basis of the theory should be carefully examined.

And so I began to investigate the details of the historical development of the theory, with the idea that I should be able to find an argument there for how the standard observational model could pertain to an essentially different physical framework that would support the metrical interpretation of cosmic expansion; as, in the meantime, I had already come across a formal mathematical result which strongly indicated to me that this should be possible.% At first, I only vaguely realised how crucial this thorough historical examination should prove to be, as it was only though carrying out the work that I came to fully appreciate the significance of historical research which James Clerk Maxwell expressed in the preface to his great \textit{Treatise on Electricity and Magnetism} \cite{Maxwell1892}, where he wrote, in connection with his own examination of the works of Michael Faraday, that `It is of great advantage to the student of any subject to read the original memoirs on that subject, for science is always most completely assimilated when it is in the nascent state'; and which Albert Einstein also noted, for a different reason, when he stated that \cite{Howard2010}
%\begin{quote}
%So many people today---and even professional scientists---seem to me like somebody who has seen thousands of trees but has never seen a forest. A knowledge of the historic and philosophical background gives that kind of independence from prejudices of his generation from which most scientists are suffering. This independence created by philosophical insight is---in my opinion---the mark of distinction between a mere artisan or specialist and a real seeker after truth. 
%\end{quote}

One of the most significant discoveries for me, which I only made in carrying out this investigation, was that the interpretation of $\Lambda$ as the essential cause of cosmic expansion is actually the one that Arthur Eddington had always fought desperately for; however, he didn't principally think of this in the way of a requisite explanation of cosmic expansion, as he had already begun cultivating the idea in the early 1920s, but actually as the most natural consequence of general relativity theory, as he considered, long before Hubble's confirmation of the cosmic expansion, that the theory could not possibly do without $\Lambda$ as a fundamental metrical constant \cite{Eddington1923}, even to the extent that he considered the prospect of dropping the cosmological constant as being tantamount to knocking the bottom out of space \cite{Eddington1933}.

To briefly sum up the first chapter of the thesis, then:---it is a detailed examination of the development and current status of modern cosmology, with special attention paid to the works which regarded $\Lambda$ in one way or another, and an eye to the development of a new theory that would be completely consistent with Eddington's interpretation of cosmic expansion.

Now, the initial mathematical result referred to above, which really paved the way for this and the rest of the analysis, came from studying the solution which describes the existence of a Schwarzschild mass in de~Sitter space, 
\begin{equation}
ds^2=-\frac{r}{r^3-{r_0}^3-(r-r_0)}dr^2+\frac{r^3-{r_0}^3-(r-r_0)}{r}dt^2+r^2(d\theta^2+\sin^2\theta{d\phi^2}),\label{SdS_pure_intro}
\end{equation}
which is reproduced here in the dimensionless coordinate system that is extensively analysed in Chapter~\ref{chap_ANC}. The Schwarzschild-de~Sitter (SdS) line-element really has two distinct forms: it describes the local gravitational field outside a non-rotating, uncharged spherical mass; but it is also a cosmological solution, in the case ${r_0}^2>4/3$, when $r>0$ is timelike (in which case the explicit form given in Eq.~(\ref{SdS_pure_intro}) directly applies). The important discovery, was that according to an observer who `naturally evolves' only according to the gravitational field described by this solution, `space' is perfectly flat, and $r$ increases with proper time, $\tau$, as
\begin{equation}
r(\tau)=(2M)^{1/3}\sinh^{2/3}\left(3\tau/2\right),\label{SdS_scalefac_intro}
\end{equation}
where $2M\equiv{r_0}-{r_0}^3$ is a constant---which is exactly the form of the scale-factor in the flat $\Lambda$CDM model.

The problem, however, is that it is not at all obvious how these coincidences could be related to the standard cosmological model, primarily because the slices of constant $r$ are not synchronous in this frame, which is tilted with the spacelike singularity at $r=0$ always at one extent of synchronous hypersurfaces, with the other spacelike direction extending towards $r=+\infty$.\footnote{It would be helpful to refer to Fig.~\ref{over-critSdS_Null}, while reading this and the following few paragraphs.} But then there seemed to be a deeper problem with this picture, because the solution actually describes \textit{all} particles at constant spatial coordinates as being `naturally at-rest'---so it does not make sense that the bundle of worldlines of all such particles would be \textit{synchronously} evolving, with particles at \textit{all} values of $r$, including those `continuously ejected from the singularity at $r=0$', comoving as they evolve through spacetime, because the gravitational potential drastically varies throughout those synchronous slices. 

In fact, the only thing that would make sense is if these particles would be comoving in the $r$-direction, along spatial slices that are tilted in this frame, so that they would always be at points of spacetime with equivalent gravitational potential.

But then it becomes important to consider particles that move through space in this coordinate system, which are therefore not `naturally at-rest'. Would these particles, as they evolve with the whole system, not remain in the same comoving hypersurface as the coherent bundle of geodesics just considered? Or would a particle that moves through space in one direction find itself existing at values of $r$ through which the rest-frame particles had already passed, and another, moving in the opposite direction, come to larger $r$ `before' everything else?

Again, the only thing that makes sense is that there would be an absolute cosmic time, described so that every particle, no matter how it would move through space, would exist on the same causally coherent hypersurface evolving along increasing $r$ from a common origin at $r=0$.

Because of the similarity of solutions, this problem then got me thinking about the validity of the standard black hole picture, as I began to wonder whether it should really be possible for an astronaut existing inside a black hole to shine two flashlights in `opposite directions', in such a way that the light from one would reach the final singularity `before' the other (consider, e.g., the evolution of null lines beyond the horizon in synchronous slices of Eddington-Finkelstein or Kruskal-Szekeres coordinates). Would this not mean that the photons that left the one flashlight would, throughout their cosmic evolution towards the final singularity, exist, in a realistic sense, \textit{later} than the astronaut, while those of the other, which were somehow shone through spacetime `towards' the past horizon, would exist \textit{earlier}? How should the notion of time that is implied by such speculation actually be justified in a causally coherent way, when $r$ is a timelike dimension, according to which things can even be said to `happen'? i.e., when, in the four-dimensional spacetime continuum, no matter how we `extend' the coordination of events, $r$ should never be thought of as a dimension through which things can be said to `move', but really as the denominator of the verb (if we should remain consistent in our interpretation of the meaning of the description)? 

Couldn't it then be instead, that regardless of the frame in which the evolution of worldlines is considered, particles travelling along every null line emanating from some given event at some $r<2m$, should progress through spacetime, along with any other particle that might be said to have existed at that same $r$ `simultaneously' with that event, so that they should also be said to exist at all later $r$, including the final singularity, `simultaneously' as well? i.e., that the existence of photons and astronauts inside a black hole could be interpreted as a comoving evolution along $r$, no matter the frame?

Similarly, we should say that there is no way of shining a flashlight in any direction of the Universe so that we could reasonably say that its light `now' exists, in four spacetime dimensions, five minutes ago, and `now' six, with timelike separation increasing, as everything also dynamically evolves in \textit{another sense}, `forward in time'. For every point in those four spacetime dimensions corresponds to an event that should take place in the Universe, so that if we want to describe instead the real motion of a particle \textit{through} something like a spacetime manifold, we need to account consistently---and analytically---for that additional notion of time.

The two facts, that the gravitational field is weak near the horizon in most black hole solutions, and that the horizon is always at $r=2m$, regardless of the coordinate system used, suggest that there really is something wrong with this general picture, in which particles are thought to continually fall through the horizon and continue on towards $r=0$ just as though $r$ were another spatial dimension---although, of course, one in which everything is `trapped' and must evolve uni-directionally---through which a particle's trajectory could be mapped, and the same path might be followed `some time later' by another particle.  For there is really no good reason to suspect, that just beyond the horizon the relativistic description should be any different from any other cosmological solution; and we should never say that a rocket flying randomly through space should accidentally find itself out of sync with the presently evolving cosmic time of a FLRW universe, just because instants of cosmic time are not synchronous in its local coordinate system and it measures proper time at a different rate; nor should we ever think of our Universe as truly evolving in cosmic time according to the standard model, while another universe, just like ours, that `big banged' a billion years later, is somehow evolving through the same four-dimensional manifold as us (as long as we think of things in terms of black hole physics, in which spacetime is like a four-dimensional arena that test-particles move through), a billion years back in the cosmic time direction; and we definitely should not think of rockets as being able to fly between the two.

And so the question, as regards the standard picture of black holes, is this: is it really correct, or even consistent, to conclude,---based on the fact that there are coordinate systems in which the trajectories of two particles can be drawn so that both fall towards $r=0$, from $r>2m$, at different `times', while a third particle `remains' outside,---that any star can have collapsed to a singularity, and more particles should continue to flood in through the horizon to the same ultimate fate, moving through the four spacetime dimensions there, in a causally coherent sense during the course of the Universe's real evolution? The answer, I believe, is no. And the resolution of this problem, which is the same as my original problem of needing a coherent description of cosmic evolution in the cosmological SdS solution, Eq.~(\ref{SdS_pure_intro}), unfortunately requires a different interpretation of the \textit{meaning} of general relativity theory---although one which is totally consistent with relativistic cosmology, and really explains why the standard model, which is consistent with both interpretations, is the most trivial possibility.

Basically, what I needed to argue, was that Einstein was already wrong when, in developing special relativity theory, he took the principle of relativity, that the laws of physics should have the same description in all frames, along with the relativistic description of simultaneous phenomena, to mean that no objective significance may be attached to any particular coordinate system, or, in another way, that the synchronous evolution of things in any frame always provides a good---ontic---noumenal description of simultaneity. 

Instead, what I needed to justify was that in all general relativistic solutions the timelike dimension actually best describes the absolute cosmic evolution of a specific three-dimensional hypersurface of spacetime, which is described just as well in all frames, according to the general covariance of the coordinate transformations---and beyond that, in a manner that is roughly consistent with the common idea of the motion of matter within a black hole's event horizon, that this real evolution occurs as a three-dimensional universe evolves dynamically \textit{through} a four-dimensional manifold, so that spacetime should not be thought of as something that really exists at all, but as the map of noumena that occur throughout that course of evolution, which can be similarly described in the proper coordinates of any reference frame, so that different noumena would be given to occur as simultaneous phenomena in different coordinate systems, but with a greater sense of simultaneity pertaining to the well-defined evolution of the universe, which would generally not be synchronous; i.e., that every event on a spacetime manifold corresponds to a hypersurface of events which truly occurred simultaneously on a different four-dimensional manifold through which that hypersurface really moves in a well-defined way, which `motion' is the same as the cosmic time of modern cosmology, as well as absolute time according to Newton's definition \cite{Newton1966}:
\begin{quotation}
Absolute, true, and mathematical time, of itself, and from its own nature, flows equably [i.e., uniformly] without relation to anything external, and by another name is called duration: relative, apparent, and common time, is some sensible and external (whether accurate or unequable) measure of duration by the means of motion, which is commonly used instead of true time; such as an hour, a day, a month, a year.
\end{quotation}

Apart from a need to sort out how general relativity could be consistently assimilated into such a cosmoogical framework---e.g.~because relativity theory is generally supposed to be at-odds with the absolute time of Newtonian physics, as indicated by the difference between the Galilean and Lorentz transformations of spacetime coordinates, or because a description of the warping of spacetime around massive particles should need to be consistently accounted for---there were other, more long-standing issues that needed to be addressed, such as Zeno's paradox, that `if everything that exists has a place, place too will have a place, and so on 
\textit{ad infinitum}' \cite{Aristotle1984}; in other words, because by adding yet another timelike dimension to the theory---a real fourth dimension other than the one of spacetime, which was now supposed to correspond to observable reality only at the present hypersurface---the theory now was, philosophically speaking, treading on dangerously thin ice.

However, I did have what I now consider to be an ill-conceived notion of time in spacetime, according to the black hole picture, to guide me; for although the `extended' coordinates of Schwarzschild spacetime remain four-dimensional coordinate systems, they still do, in the way that their causal description has been interpreted, support such a five-dimensional idea; i.e., because they describe a four-dimensional set, $\{r,t,\theta,\phi\}$, bounded in $r$ ($<2m$) and spherically closed, which \textit{evolves} in such a way that particles can be \textit{continually dropped into it}. So although such coordinations of the Schwarzschild metric do not explicitly add another dimension, but only `extend' these ones, it is, I believe, through mis-interpretation of the causal evolution that is described in the new systems, according to that same aspect of Einsteinian relativity that I needed to argue against, that this fifth dimension is anyhow realised.

At this point, the project had really become one on the philosophy of Time, which is a very old problem, and so the rigorous justification I needed to give for the extra-dimensional type of `presentist' theory I wished to propose required analysis of the basic developments that were made on both sides of the argument, throughout history. The following is a very brief synoptical account of that:---

Twenty-five hundred years ago, Parmenides of Elea had a thesis: all of space and all of time simply exist as a singular (spherical) block universe, and the perceptions of duration, free will, and change are purely illusory. Actually, that statement is not even correct, because Parmenides thought that space and time couldn't actually exist at all, but that all that exists is substantive matter, which therefore exists as one four-dimensional \textit{plenum}. Parmenides worked out this theory deductively, from a single axiom: What-is is all there Is.

In response to Parmenides, the Atomists, Leucippus and Democritus, proposed a new axiomatic physical principle: What-is-not, which they called the `void', is every bit as Real as What-is, which is atoms of all shapes and sizes, and the atoms move about in the void, bringing about change as they do so. 

But the Atomists' void was so at-odds with the common sense-based notion that there truly is substance everywhere, with air or liquid filling even the tiniest pores of every otherwise solid object, that even Aristotle, who sought the same goal of breaking through the barrier that Parmenides had set up, and thereby reconciling his analytic physical theory with fundamental change in reality, in accordance with the Heraclitean thesis of time as the uniform flux of all things, chose instead to admit the requirement of a supernatural prime mover who would achieve the same end as the Atomists' void---as he did allow for the reality of `place', but developed his theory so that this `place' would be continuously full of matter, which \textit{rearranged} itself in time, according to a well-ordered duration that was achieved through the abstract metaphysical prime mover.

Centuries later, Newton developed his theory as one that was very much in line with the principles of ancient atomism, which did not admit the presence of a supernatural being anywhere in physical reality except in the initial set-up of the solar system. According to Newton's theory, things moved through absolute space and time which are ontologically as real as the substantial material that exists in them, as the whole system of physical reality would equably endure, with the same order that Aristotle's prime mover was intended to achieve, which requirement was probably made the most clear after Aristotle, through St.~Augustine's rationalisation. 

It is my opinion that Einstein, who, e.g.~along with Ernst Mach and Ren\'{e} Descartes, was more inclined towards the Parmenidean way of thought when it came to the age-old debate on the independent reality of a substantive void, than, say, the Atomists and Newton, actually \textit{developed} or \textit{cultivated} his interpretation of the meaning of relativity theory so that it would be in line with that position, which always inevitably has led through, deductively, to the requirement of a block universe. 

This is evidenced, e.g., in `Relativity and the Problem of Space', the fifth appendix to \textit{Relativity: The Special and General Theory}---e.g., where he argued that \cite{Einstein2001}
\begin{quotation}
\ldots. Mach, in the nineteenth century, was the only one who thought seriously of an elimination of the concept of space, in that he sought to replace it by the notion of the totality of the instantaneous distances between all material points. \ldots

[In Newtonian mechanics], ``physical reality'', thought of as being independent of the subjects experiencing it, was conceived as consisting, at least in principle, of space and time on one hand, and of permanently existing material points, moving with respect to space and time, on the other. The idea of the independent existence of space and time can be expressed drastically in this way: If matter were to disappear, space and time alone would remain behind (as a kind of stage for physical happening).

\ldots 

We are now in a position to see how far the transition to the general theory of relativity modifies the concept of space. In accordance with classical mechanics and according to the special theory of relativity, space (space-time) has an existence independent of matter or field. In order to be able to describe at all that which fills up space and is dependent on the coordinates, space-time or the inertial system with its metrical properties must be thought of at once as existing, for otherwise the description of ``that which fills up space'' would have no meaning.\symbolfootnote[1]{If we consider that which fills space (e.g.~the field) to be removed, there still remains the metric space in accordance with 
\begin{equation}
ds^2={dx_1}^2+{dx_2}^2+{dx_3}^2-{dx_4}^2,\label{M4}
\end{equation}
which would also determine the inertial behaviour of a test body introduced into it. [Einstein's footnote].} On the basis of the general theory of relativity, on the other hand, space as opposed to ``what fills space'', which is dependent on the co-ordinates, has no separate existence. Thus a pure gravitational field might have been described in terms of the $g_{ik}$ (as functions of the coordinates), by solution of the gravitational equations. If we imagine the gravitational field, i.e.~the functions $g_{ik}$, to be removed, there does not remain a space of the type [Eq.~(\ref{M4})], but absolutely \textit{nothing}, and also no ``topological space''. For the functions $g_{ik}$ describe not only the field, but at the same time also the topological and metrical structural properties of the manifold. A space of the type [Eq.~(\ref{M4})], judged from the standpoint of the general theory of relativity, is not a space without field, but a special case of the $g_{ik}$ field, for which---for the coordinate system used, which itself has no objective significance---the functions $g_{ik}$ have values that do not depend on the co-ordinates. There is no such thing as an empty space, i.e.~a space without field. Space-time does not claim existence on its own, but only as a structural quality of the field.

Thus Descartes was not so far from the truth when he believed he must exclude the existence of an empty space. The notion indeed appears absurd, as long as physical reality is seen exclusively in ponderable bodies. It requires the idea of the field as the representative of reality, in combination with the general principle of relativity, to show the true kernel of Descartes' idea; there exists no space ``empty of field''.
\end{quotation}

Therefore, apart from the fact that he really was required by developments of modern physics to admit that there is vacuous space between particles of matter, Einstein really had come back as close as he could to Parmenides' theory of one block universe. In fact, the relativistic description of a block universe was first expounded by Hermann Minkowski, when he invented `spacetime' explicitly as an absolute four-dimensional world in 1908 \cite{Einstein1952}. 

In 1966, rigorous proofs were given by Wim Rietdijk \cite{Rietdijk1966} and Hilary Putnam \cite{Putnam1967} which confirmed that the standard interpretation of special relativity theory due to Einstein, that every inertial observer is objectively as justified in claiming that the simultaneous phenomena they observe in their frames are really simultaneous noumena, really does require a block universe. Actually, this is obvious: if we integrate all possible simultaneities of all possibly existing simultaneous special relativistic observers, even at one particular instant in a specified frame, we find that the four dimensions of Minkowski spacetime must absolutely exist simultaneously, so that we can `imagine that everywhere and everywhen there is something perceptible', as Minkowski put it \cite{Einstein1952}.

The Rietdijk-Putnam determinacy proofs were arguably strengthened later by Howard Stein's \cite{Stein1968,Stein1991} attempt to refute them \cite{Callender2000,Savitt2008}. For, by arguing that special relativistic spacetime will be indeterminate if, and only if,\footnote{The argument only works if spacetime is all that really exists, and no more structure is added to it, as it exists in this indeterminate state.} every event outside any observer's causal past is, in a general or absolute sense corresponding to each observer individually, indeterminate here-now, Stein effectively showed that any such `light-cone argument' in favour of indeterminacy should fail if it is possible to say that even one event that lies outside our causal past, such as the emission of a photon towards the Earth, from \textit{somewhere} on the Sun's photosphere, \textit{sometime} in the past eight minutes, should have occurred without our exact knowledge of the where and when of it. Thus, e.g., by positing that `at least one event in the universe shares its present with another event's present', which he considers to be `the thinnest requirement one could put on becoming', Craig Callender transforms `Stein's ``possibility'' theorem into a ``no go'' theorem for objective becoming in Minkowski spacetime' \cite{Callender2000}.

Stein's theory basically corresponds the description of Minkowski spacetime given by Eddington in \textit{The Nature of the Physical World} \cite{Eddington1929}. In my opinion, such a radically phenomenalist standpoint, which goes contrary to all reason just because it can, according to the principle of scientific verifiability, is metaphysically useless. Even so, it is worth briefly considering whether Stein's proposal for reconciling becoming with relativity theory should have anything to do with standard physics.

To begin with, we might contrast Stein's idea with the basic justification for the standard black hole picture; for then it should be immediately obvious how inconsistent it would turn out to be, if any physicist would claim that they seriously subscribe to this view rather than the Einsteinian one---as contemporary wisdom has it that black holes presently exist all throughout the universe, even though it is impossible that any event that would have occurred at a timelike radius inside a black hole's event horizon, or even precisely at the horizon, could ever find its way into the past light-cone of an external observer. In fact, according to Kip Thorne \cite{Thorne1994}, the basic result that convinced the global community of theoretical physicsts that black holes should presently exist in the Universe, was David Finkelstein's  \cite{Finkelstein1958} discovery of a coordinate system for the Schwarzschild metric, that describes the synchronous evolution of an observer who remains outside, along with a particle that passes through the horizon and on to the central singularity, in a finite interval of coordinate time. 

However, according to Stein, all points outside the past light-cone of here-now mustn't yet have come to be if relativity theory is to describe anything but a four-dimensional block universe; so it cannot be concluded that any particle anywhere in the entire Universe has ever dynamically reached the event horizon of a black hole, because all such events should necessarily have to fall outside the causal past of every external observer---i.e., there is no way of saying definitely that the proper time of any particle that is taking part in gravitational collapse has ticked past whatever finite interval it should take to get arbitrarily close to the horizon. The interpretation of Finkelstein's result as providing evidence that the process of collapse will already have occurred therefore requires the assumption that the description of synchronous space has objective meaning in that frame---and thus, probably every other, because there is absolutely no justification for assuming that Eddington-Finkelstein synchronicity should be somehow special,---which in turn demands a block universe if one actually cares to be honest, along with Einstein,\footnote{Thus, he admitted that `To us believing physicists the distinction between past, present, and future has only the significance of a stubborn illusion' \cite{Foelsing1997}.} about the formal implications of one's principal interpretation of the theory. 

Another example of the fact that relativists have, in principle---if not rigorously in subsequent application,---commonly sided with Einstein et al.~rather than Stein, when interpreting the meaning of relativity, is Roger Penrose's treatment of the relativity of simultaneity in \textit{The Emperor's New Mind} \cite{Penrose1989}, where he presents a similar argument to the one given by Rietdijk and Putnam, noting the same implication for determinism. In fact, Penrose even goes so far as to claim, in an endnote to his statement that `The ``now'' according to one observer would not agree with that for another', that `Some relativity ``purists'' might prefer to use the observers' light cones, rather than their simultaneous spaces. However, this makes no difference at all to the conclusions.' In fact, it makes no such difference \textit{except} in the extreme case proposed by Stein---which, we must therefore infer, Penrose did not register as a realistic possibility. 

Therefore, although we do commonly think of general relativity as a dynamical theory, which would describe the continually evolving local interaction of matter with a four-dimensional metric space that it dynamically warps, the original interpretation of the meaning of the general relativity of simultaneity, due to Einstein, which is still commonly assumed by physicists, leads deductively through to a block universe description, according to which the supposed dynamism is not fundamental.

As Robert Geroch put it \cite{Geroch1978}, `There is no dynamics within space-time itself: nothing ever moves therein; nothing happens; nothing changes. \ldots one does not think of particles ``moving through'' space-time, or as `following along'' their world-lines. Rather, particles are just ``in'' space-time, once and for all, and the world-line represents, all at once, the complete life history of the particle.'\footnote{cf.~my above remark, that the timelike dimension of spacetime does not describe something through which anything actually moves, but is simply the denominator of the verb.} Or, as he finally concluded, `[Perhaps the distinctive feature of Einstein's theory of relativity] is the introduction of the idea of space-time---the idea that everything of physical interest in the world, including oneself, is to be set down therein once and for all. In the Galilean view, space-time is a luxury; in relativity, a necessity.'

This direct consequence of the standard Einsteinian interpretation of the theory---which implies that there is absolutely no purpose in undertaking any more physical investigations that might lead us to a better understanding of the Universe, since all of it, and everything we do in our lives, should be strictly determined anyway, all of it substantively existing as a fixed four-dimensional block---is not only the most significant, but also the most commonly overlooked problem of modern physics; and it should be impractical, if not impossible, to try to resolve any of the other problems connected with relativity theory so long as it is inconsistently treated, within the standard framework, as a fundamentally dynamical theory, and thereby used to describe pseudo-dynamical processes that are essentially paradoxical.

But as Stein correctly pointed out, first in \cite{Stein1968}, in relation to Putnam's argument, but also in \cite{Stein1969}, in relation to Parmenides', `An argument that leads to a truly paradoxical conclusion is always open (if it escapes conviction for fallacy) to construction as a \textit{reductio ad absurdum}.' This is precisely what the ancient Atomists did, in order to reconcile a physical theory of motion, according to the (otherwise arguably gratuitous) premiss that space really does have an independent ontological existence. For the case of special relativity theory, it is worthwhile to reproduce Stein's more detailed critique of Putnam's argument as well \cite{Stein1968}:
\begin{quotation}
\ldots Putnam's conclusion---that all things, for all time, are real---contradicts the ``man on the street's view.'' What sort of argument is it, then, that contradicts its own premise?

From the strictly deductive point of view, such an argument is a \textit{reductio ad absurdum}, i.e., a refutation of the conjunction of its premises; and if we consider the ``man on the street's view'' to be the only questionable premise of the argument, the conclusion is that the man on the street is in error \ldots. But an argument whose conclusion is incompatible with its premises may be viewed in another way: as a heuristic, rather than a strictly deductive, argument. The major difference is that whereas the formal ``conclusion'' of a \textit{reductio ad absurdum} is of no interest in itself (as being ``absurd''), a heuristic argument can lead to the acceptance of a positive conclusion which, in some sense, \textit{corrects} its premises. The principles of such argumentation are notoriously hard to codify; and it seems to me that, just for this reason, if one regards the process as a serious one---and not just a frivolous exercise or game---, one is obliged in pursuing it to seek the greatest possible clarity at all stages.
\end{quotation}

The reason for including this lengthy quotation here, is that it touches on a point in the epistemological argument that I've given in Chapter~\ref{chap_AER} (esp.~\S~\ref{sec_ANE}). For, the theories of Einstein and Parmenides really must both be incorrect for the same basic reason: they work from axioms that prove to be untenable because they lead to an absurdity. Therefore, it is the axiomatic basis of Einsteinian relativity that needs to be examined critically in order to reduce the theory to a new one which is not absurd.

Now, although Stein should therefore be entirely correct in his assessment of how the error comes into relativity theory, I believe he has made an error in pinning down what that error is. For there really should be nothing wrong with the `man on the street's view', so long as he stands motionless with respect to all of his surroundings. The problem, I believe, occurs when we try to attach the same significance to the view of the `man in the train', or the `man in the cabin of the boat', who, although they are also inertial observers, cannot really be thought of as being in a true state of rest, any more than we think the Earth is, when we interpret the meaning of the dipole anisotropy in the cosmic microwave background (CMB). The fact that we now possess such \textit{strong empirical evidence in favour of the existence of a true state of absolute rest}, as the CMB provides, suggests that such pre-cosmological arguments for the significance of the relativity of inertia, as Einstein employed in motivating his interpretation of the relativity of simultaneity, can no longer be given the weight that they once were. Instead, we should utilise---to its fullest extent---the fact that the empirical evidence suggests that there actually is an absolute frame of rest, as Hermann Weyl already realised in the early 1920s \cite{Weyl1923,Weyl2009,Weyl1924,Weyl1930,Robertson1933,Bondi1961,Ehlers2009,Rugh2010}, and aim for an interpretation of relativity theory that would describe that state consistently in local frames as well.

Accordingly, the Atomist-Newtonian-type interpretation of relativity theory that is described in Chapter~\ref{chap_DIRT} begins from the axiomatic standpoint there really is a four-dimensional void through which a three-dimensional material universe evolves with well-defined absolute cosmic time. I mean, that absolute cosmic time is a common property of every particle that makes up the universe, so that the present association of particles that may also move randomly with respect to each other, through the universe in cosmic time, evolves uniformly along the four-dimensional void.

This can be reconciled with special relativity theory according to the following kinematical principle: as the universe moves through the void, which in this case is four-dimensional Euclidean space, no particle in the equably evolving universe, which is always a well-defined hyperplane, can ever move further through the universe than it does in the direction of cosmic evolution. The map of occurrences in that well-defined three-dimensional universe is an evolving block spacetime that might somehow have been etched into the void, actually remaining in the past, as in the evolving block universe theory which has recently been discussed by George Ellis \cite{Ellis2007,Ellis2008,Ellis2010} (although there is no void in that framework); however, the past block of spacetime events really doesn't exist any more than the future one in the void/absolute cosmic time interpretation, as all that exists is the four-dimensional void and the present universe evolving in absolute cosmic time, carrying information about past occurrences along with it, e.g.~in the form of photons which move at the cosmic speed limit.

The absolute cosmic time and speed limit are so far assumed axiomatically, as physical laws; and, in relation to the standard interpretation of special relativistic spacetime, the same properties of the physical description of spacetime events that did occur, can be subsequently recovered if (in the Euclidean space on which those events should be mapped) we simply replace the coordinate that describes the direction of flow of cosmic time with a pure imaginary number, so that what is described is an evolving block Minkowski spacetime, with its Lorentzian signature and everything that entails. 

This interpretation of the physical description is similar to the one that was initially made by Hendrik Lorentz \cite{Lorentz1904,Einstein1952} and Henri Poincar\'{e} \cite{Poincare1913}---which was, e.g., advanced later by John Bell \cite{Bell1976,Bell1987};\footnote{A recent volume, dedicated to the exposition of this alternate framework, is \cite{Craig2008}.} however, some nontrivial structure has been added to the description of physical reality, which further clarifies this position as one which is in fact consistent with the standard description of relativistic covariance, according to which the phenomena of time dilation and length contraction result purely from the finite speed of light through the universe and its invariant local measure. For we can now elaborate on the notion that the simultaneous events occurring as `local time' progresses in an arbitrary inertial frame, do not necessarily \textit{really} occur simultaneously, viz.~in Lorentzian `true time' \cite{Lorentz1904,Einstein1952}, by noting that at any instant in the frame of a relatively moving observer, the events that really do occur simultaneously (simultaneous noumena) take place on a spacelike hyperplane that extends to the `past' and `future' in the local coordination of four-dimensional \textit{reality} at any instant, and, similarly, that the synchronous hyperplane extends into the void, where no universal noumena are occurring; i.e., there is an objectively well-defined distinction between the simultaneous phenomena that the theory describes as events in the four-dimensional spacetime continuum, which are coordinated differently by observers in relative motion, and those which are in fact simultaneous noumena. Then, as the special relativistic universe is clearly and objectively well-defined in all frames, it will be described, in the frame of an inertial observer who is not in a state of absolute rest, as evolving at some nonzero angle with respect to the plane of synchronicity, although the worldlines of all particles that travel through the universe at rates up to the cosmic speed limit still must be timelike, and the null worldlines describing the paths of photons will remain phenomenally invariant. Therefore, depending on the coordinate system that is used to describe the noumenal progression of events, the phenomenal measures of true lengths and durations will also transform covariantly.

Following Einstein, one might be tempted to label the absolute time and void of this theory as superfluous, since the description of phenomena is the same without them; but, lacking these additional elements, Einsteinian relativity theory is fundamentally at-odds with the common sense perceptions of Time and free will, and is therefore absurd. It should therefore be noted that this Lorentz-Poincar\'{e} interpretation of the theory is more objective, because, along with everything that is accounted for similarly as in Einstein's limited theory, it is essentially consistent with these most common facts of observation as well.

Another objection that might be made, is to the artificial introduction of Lorentzian signature into the special relativistic spacetime metric; however, this simply results from subsequently defining the local time axis as the worldline of an inertial observer who remains at rest in that frame, in conjunction with the requirement that the speed of light should be measured to have the same constant value in all such frames; i.e., by assuming the standard inertial and causal structures of the theory. The physics is thus already contained in the kinematical principles, while the metrical structure defined through the Wick rotation simply makes that easier to follow, i.e. through Minkowski's spacetime formalism. Therefore, that definition cannot be considered any more or less artificial than the Lorentzian signature that is generally assumed in relativity theory.

But because the Lorentzian signature of the metric may in fact be fundamental, as it would be if general covariance and the constant finite speed of light did really depend on it, the more important point to consider, is whether there is a more realistic way of formulating the theory, so that the Lorentzian signature should be described, instead, as a basic property of physical reality. 

Together with this problem, it was then also important to consider whether the void itself should truly be dynamically warped by gravitational mass, just as spacetime is warped in frames describing the existences of massive bodies. The idea that I had to begin with, according to the cosmological SdS solution, was that the void should indeed be warped by a massive universe evloving along $r$ from an initial singularity, with the curvature of the external (future) void given according to the Jebsen-Birkhoff theorem \cite{Jebsen1921,Jebsen2005,Birkhoff1923}, so that the timelike nature of $r$ would be forever induced according to the field theoretic description of the massive universe's interaction with the de~Sitter void in which it was embedded.\footnote{The non-locality of this description seemed to be no less troubling than that of the standard black hole picture, in which the curvature of spacetime throughout $r<2m$ should be induced by a future spacelike singularity that was already created in the principal collapse, to which all particles that subsequently fall into the hole are teleologically attracted.}

However, the eventual theory that I came to, although effectively equivalent to this idea, is essentially different.

The problem that I encountered right away, was that a timelike direction, in a metric with Lorentzian signature, is probably required to begin with, in order to describe the warping of spacetime around an existing mass. The reason (according to a well-known calculation) is simple: if one tries to describe the instantaneous curvature of the four-dimensional void around an ultimate point mass, assuming that it must be isotropic, the solution of Einstein's equation is simply the metric for maximally symmetric space. 

This can scarcely be called a proof that gravitational mass shouldn't somehow actively warp the basic fabric of physical reality; and no doubt such a theory could be formalised some other way---i.e., some way that would not also require the block existence of the four-dimensional field in conjunction with that space, as in Ellis' evolving block universe theory, in which the past-and-present field that exists \textit{is} the basic fabric of reality, which is warped around massive bodies as general relativity theory is supposed to describe;\footnote{However, it should be noted that the only reason this would be possible, is that Ellis has forsaken the Einsteinian interpretation of the relativity of simultaneity as well, and assumed an absolute time so that the synchronous present in all frames is not also the noumenal present---for otherwise Einstein's argument in \cite{Einstein2001} proves a \textit{reductio ad absurdum} in the same spirit as Putnam's: viz., that the standard idea that the spacetime field should really dynamically mould its accompanying space, taken in conjunction with Einstein's interpretation of the diffeomorphism invariance of that field, ends up requiring that the assumed dynamical warping of space doesn't actually take place in a fundamentally evolving way, but is moulded into a four-dimensional static block.} but, under the assumption that in reality there \textit{is} a four-dimensional space with an independent ontological existence, there is also no good reason to \textit{require}, just because our theory describes the progressively occurring spacetime field as being dynamically warped around gravitational mass, that this should be the case for the background void as well---especially when our theory also tells us other things, like the facts that \textit{all} particles that exist must have energy (canonical conjugacy of $E$ and $t$), that mass is equivalent to energy, that there is no effective difference between gravity and inertia, and that energy is the timelike component of the four-momentum.

When I realised that such points as these---combined with the fact that the supposed dynamical warping of the void around a gravitational mass is not naturally described through the reduced principles of the theory I was now attempting to synthesise---seemed to indicate that mass might actually be described as corresponding to the inertial path of a particle through the void, the answer to the question of whether there should be a basic Lorentzian signature in physical reality came very quickly, according to the strong empirical evidence suggesting a basic Cartesian ordering in Nature---since, according to this non-trivial fact, the void, now understood to be an absolute space according to Newton's definition,\footnote{sc.~\cite{Newton1966}, `Absolute space, in its own nature, without relation to anything external, remains always similar and immovable. Relative space is some movable dimension or measure of the absolute spaces; which our senses determine by its position to bodies; and which is commonly taken for immovable space; such is the dimension of a subterraneous, an aerial, or celestial space, determined by its position in respect of the earth. Absolute and relative space are the same in figure and magnitude; but they do not remain always numerically the same. For if the earth, for instance, moves, a space of our air, which relatively and in respect of the earth remains always the same, will at one time be one part of the absolute space into which the air passes; at another time it will be another part of the same, and so, absolutely understood, it will be continually changed.'} should really have to be maximally symmetric.

Using Cartesian coordinates, these few spaces are straightforward to classify (see Eqs.~(\ref{R5}) -- (\ref{eigenvalue}), and the surrounding discussion): in $n$ dimensions, they are either $n$-dimensional flat space, or $n$-dimensional surfaces with nonzero constant sectional curvature, which may be described as hyperspheres with constant radius of curvature, embedded in $(n+1)$-dimensional flat space; so, if we begin by assuming that the coordinates describing the maximally symmetric space should all be real, we find, e.g., that the flat, maximally symmetric real space must be Euclidean space, which correspondence to special relativity theory ammounts to a re-definition of one coordinate as purely imaginary, as discussed above. If the $(n+1)$-dimensional flat embedding space is Euclidean, the only maximally symmetric hypersurface is the closed $n$-sphere. If, on the other hand, the embedding space is Minkowski space (by virtue of the embedded sphere that is described; see Eq.~(\ref{sphere2})), then the maximally symmetric space can have negative curvature, as described by the open $n$-sphere with imaginary radius of curvature (hyperboloid of two sheets), or positive curvature, as described by the open $n$-sphere with positive radius of curvature (hyperboloid of one sheet). The latter is de~Sitter space, which is a real manifold with one timelike tangent vector at every point, so that its metric has Lorentzian signature; and the former, which is anti-de~Sitter space, has a purely spacelike basis, so its metric is positive-definite.

The proof of this coherent description, which is given in \S~\ref{sec_RPT}, is significant for a few reasons. For one thing, it is commonly said that the open hyperbolic slice of Minkowski space with negative curvature, is the negative curvature analogue of the closed sphere---i.e., that the geometries described by
\begin{equation}
ds^2=\frac{dr^2}{1-Kr^2}+r^2{d\Omega_{n-1}}^2
\end{equation}
are (closed) spherical or (open) hyperbolic, depending on whether the curvature constant $K$ is positive or negative, respectively. For basic topological reasons, such a correlation actually should not be drawn.

For example, when $K>0$, there really is no good reason for assuming that the coordinate $r$ goes only from 0 to $K$, and that the space that the line-element \textit{really} describes is the closed sphere which continues on in a regular manner below the equator, as described, e.g., in \cite{Weinberg1972}; for it is just as true that for $1/\sqrt{K}\leq{r}\leq\infty$, the line-element describes de~Sitter space---just as the curve $y=\sqrt{1-x^2}$ is hyperbolic, with $y$ purely imaginary, for $x^2\in\mathbb{R}\geq1$; i.e., it is a one-dimensional surface with real domain, embedded in the complex plane.

And, furthermore, if we should be daunted by the complexity of this embedding space, then we must be so consistently, and say that we don't like the curve $y=\sqrt{-1-x^2}$, $x\in\mathbb{R}$ either. But we don't do this:---we describe this very curve---or its higher-dimensional counterpart, given by rotations about the (imaginary) $y$-axis---as open hyperbolic space with negative curvature; and we call it the analogue of closed spherical space with positive curvature; and we describe de~Sitter space as a maximally symmetric pseudo-Riemannian manifold, with Lorentzian signature and positive curvature.

However, the maximally symmetric space with negative curvature, or disconnected open sphere (hyperboloid of two sheets) with imaginary radius of curvature embedded in (induced) Minkowski space of one higher dimension, is not the analogue of the closed sphere with positive radius of curvature, but of de~Sitter space, which may be described as a \textit{real} maximally symmetric space with positive curvature, or a connected open sphere (hyperboloid of one sheet) with positive radius of curvature embedded in Minkowski space of one higher dimension. And the latter is every bit as \textit{Riemannian} as the former. 

And, for what it's worth, it may also be noted that the closed sphere should not necessarily be described as a surface embedded in Euclidean space, but may also be embedded in de~Sitter or anti-de~Sitter space as well; i.e., spherical space may be described as a maximally symmetric closed hypersurface of any one of these three open space forms.

And according to this coherent description, since the basic metric tensor of anti-de~Sitter space is positive-definite, it should be noted that any description of an anti-de~Sitter spacetime requires one coordinate of the basic metric to be defined as imaginary, in the same \textit{ad hoc} manner as is done in constructing special relativistic spacetime.

Now, this is really all very significant because of the fact that Einstein continually claimed that general relativity theory would be logically simplest if $\Lambda\equiv0$. However, if we consider that the Lorentzian signature of spacetime in relativity theory should be related to a basic property of Nature, so that Physical Reality would \textit{essentially} give rise to a generally covariant description, then it is really logically the simplest---whether or not we should admit the existence of a void into our physical theory, or whether such a void, if admitted, should be deformed by any absolute rest-mass that moves through it, or should itself be an absolute space; i.e., in any case---to assume a positive cosmological constant in Einstein's equation, so that its most basic solution should describe a real four-dimensional Riemannian manifold with intrinsic Lorentzian signature.

In fact, it is important here to briefly describe, as best I can in a few short words, one (arguably the only) negative aspect of Einstein's epistemological method. In order to obtain the greatest logical simplicity, he felt that one should begin with a physical theory, and sort out how it may have been derived according to the fewest number of axioms (or mathematical terms); and then cast aside, as superfluous, all others that might initially have been thought to be realistically assumed, before proceeding to refine the theory in order to bring it in line with even more facts of observation, taking care to admit only those new principles or axioms which seem to be the most reasonable; then, with such a reasonable and logically simple theory in hand, he should work to develop an acceptable level of comfort when certain aspects of the theory don't naturally make sense, as paradoxical inferences may be encountered.

But there is a fault in this, as he did search principally for logical simplicity in the axioms of his physical theory, even at the sake of some loss of intuitive clarity at the fundamental level: for while the historical development of physical theory does strongly suggest that the correct theory should be logically the simplest, no one is fit to judge, while speculating about what the principal axioms of that theory should be, which ones will give rise to the logically simplest theory that may subsequently be deduced.

For example, while he did not immediately realise that his interpretation of the various coordinate representations in special relativity theory should amount to the description of a block universe, he did eventually realise and accept this fact, which, although there is no definitive means of proving that it should not truly be, still seems extremely unrealistic. Another paradoxical aspect of Einstein's theory is noted in his autobiography, where he wrote that in developing general relativity theory, after he had realised the principle of general covariance in 1908, he could not complete the thoery until he had spent seven difficult years convincing himself that the coordinates should have no immediate metrical meaning \cite{Einstein1951}. The same level of comfort with this paradoxical notion is carefully developed, e.g., at the beginning of Charles Misner, Kip Thorne, and John Wheeler's \textit{Gravitation} \cite{Misner1973}.

In contrast, the coordinates do have immediate metrical meaning in the theory developed in this thesis. And this begins by framing the de~Sitter manifold with a particular Cartesian coordinate basis, which sets it up as the absolute background space through which the three-dimensional universe should really evolve, and the ideal evolving block spacetime should emerge. Then there is only one realistic option for the description of a three-dimensional universe that evolves in a causally coherent manner: the comoving, torsionless closed 3-sphere which evolves in the timelike direction of the corresponding five-dimensional Minkowski embedding space. This is contrasted with more usual coordinations of the de~Sitter sphere---i.e., the statical Eddington frame \cite{Eddington1923} and the comoving Lema\^{i}tre-Robertson frame \cite{Lemaitre1925b,Robertson1928}---towards the end of \S~\ref{sec_RPT}, in an analysis which clearly illustrates the illegitimacy of those descriptions as realistic cosmological models, so long as metrical significance of the coordinates is granted and de~Sitter space is not conceived as an abstract manifold.

Now, an interesting aspect of this universe, as opposed to the universe of special relativity theory, is that the only inertial particles---i.e., which move with constant velocity through it, as cosmic time progresses---whose paths are actually geodesics of de~Sitter space, are those which remain perfectly at rest. All other inertial particles are accelerated through the universe. This was already realised by Henri Bacry and Jean-Marc L\'{e}vy-Leblond in 1968 \cite{Bacry1968}, who named this particular Lie group the \textit{Newton group}, but failed to realise that the accelerations of these inertial particles might, in their proper frames, be perceived as mass.

It is worthwhile to briefly discuss the content of Bacry and L\'{e}vy-Leblond's paper, because, along with the fact that it corroborates the above argument, that $\Lambda>0$ is really the logically simplest possibility for general relativity theory, the formalism that they develop will be useful for outlining the rest of the thesis in a clear and concise way.

First and foremost, what Bacry and L\'{e}vy-Leblond proved \cite{Bacry1968}, was that there are eight different types of Lie algebras for kinematical groups that one can reasonably propose, and that the \textit{de~Sitter group} corresponds to the most general of these, so that all the others can be recovered as limiting cases of it, according to an In\"{o}n\"{u}-Wigner `contraction scheme' \cite{Inonu1953}. The article was highlighted a few years later by Freeman Dyson \cite{Dyson1972}, who noted, e.g., that the de~Sitter group is the only simple group among these.

The first of the two types of contraction which are relevant here, is the so-called `speed-space' contraction, whereby the generators of space-translations and inertial transformations are considered in their infinitesimal limits. This contracts the relative time Lie algebras to the absolute time Lie algebras; e.g., the speed-space contracted de~Sitter group is the Newton group, which is relevant for us, and will be discussed further shortly.

The other interesting type of contraction is the `speed-time' contraction, in which the time-translation generator and inertial transformation generators are sent to zero. Just as the relative time Lie algebras are contracted to the absolute time Lie algebras through speed-space contractions, so the relative space Lie algebras are contracted to absolute space Lie algebras through speed-time contractions---which describe spatial slices, in some particular proper inertial frame, as being similar for all times. Bacry and L\'{e}vy-Leblond wrote that this process `only yields groups describing intervals connecting events without any causal connection, hence without much physical applications' \cite{Bacry1968}.

However, I think this was too hasty a conclusion, since the contraction doesn't necessarily need to be taken to mean that \textit{nothing} can move through space (i.e., that we have to set $c\rightarrow0$ and thereby refuse to describe the motions of test-particles), but merely that space remains absolute relative to \textit{something} that is always described as being at rest;\footnote{Analogously, in the case of speed-space contractions, we don't actually have to say that the absolute time that they define is also the proper time of \textit{all} observers, and thus the same time-coordinate of all frames. This is why we can still reasonably define an absolute time in special relativity theory, contracting Einstein's \textit{Poincar\'{e} group} representation to the Lorentz-Poincar\'{e} \textit{Galilei group} representation of the theory.} in fact, this doesn't even need to mean that that something is \textit{actually} at rest, in an absolute sense. However, it is relevant to note that such a description of space cannot pertain to a situation in which its geometry is generated by a local action causally emanating from a massive body.

In relativity theory, one very important group to which this contraction pertains, was already noted in a way by Eddington in 1923 \cite{Eddington1923}, when he sought to describe the relatively symmetrical properties of an ultimate particle of matter, about which photons could be described to move, in the `radial' direction, `forever' at the same speed---i.e., so that the surrounding spacetime description could be locally restricted in algebraic form, to
\begin{equation}
ds^2=-U(r)dr^2-V(r)(r^2d\theta^2+r^2\sin^2\theta{d\phi}^2)+W(r)dt^2.\label{Edd_contraction_intro}
\end{equation}
This is Eddington's equation (38.12), as he referenced it in \S~66 of \cite{Eddington1923}. In fact, because this corresponds to a speed-time contraction, $U$, $V$, and $W$ should most generally be written as functions of $t$ as well, since the line-element should describe the non-causal curvature of space surrounding a particle at rest at some time, and \textit{similarly} for all times; however, that potential $t$-variability can be gauged away when the metric is required to be a solution of Einstein's equation, according to the Jebsen-Birkhoff theorem. 

In Chapter~\ref{chap_ANC}, the local SdS solutions to Einstein's equation are analysed according to the interpretation that the line-element corresponds to such a speed-time contraction of the de~Sitter spacetime field surrounding an ultimate particle of matter moving with arbitrary inertial velocity through a framed de~Sitter manifold. More will be said about that analysis shortly, which should be the most clear after the relevant speed-space contraction has been discussed.

By sending the space-translation and inertial transformation generators to zero, speed-space contractions distinguish a particular bundle of worldlines of particles that coherently evolve \textit{as} a (three-dimensional) universe with absolute time, which therefore satisfies Weyl's postulate. Such contractions of the de~Sitter group are the easiest to conceptualise if it is thought of as a framed manifold:---the four-dimensional hyperboloid of one sheet, embedded in five-dimensional Minkowski space. The trivial contraction, which is analogous to the pure de~Sitter form of Eq.~(\ref{Edd_contraction_intro}), distinguishes the inertial \textit{geodesics}, which actually do remain at rest in the framed manifold, as the fundamental worldlines. The corresponding Robertson-Walker (RW) metric, describing the evolution of this universe---i.e., describing this Newton group---from this perspective, is
\begin{equation}
ds^2=-dt^2+\frac{3}{\Lambda}\cosh^2\left(\sqrt{\frac{\Lambda}{3}}t\right){d\Omega_3}^2,\label{dS_comoving_univ_intro}
\end{equation}
where $d\Omega_3$ describes the cosmic 3-sphere which contracts until $t=0$, and expands thereafter.

However, since it has already been noted that these worldlines describe the only non-accelerated inertial paths in the universe, it seems most reasonable to suppose that a particle following such a path should be massless; indeed, according to the statical Eddington line-element, given by the speed-time contraction around such a geodesic, a particle that exists as such, remaining forever at $r=0$ in
\begin{equation}
ds^2=-\left(1-\frac{\Lambda}{3}r^2\right)c^2dt^2+\frac{dr^2}{1-\frac{\Lambda}{3}r^2}+r^2(d\theta^2+\sin^2\theta{d\phi^2}),\label{deSitter_statical_intro}
\end{equation}
should be massless. 

Therefore, the proposed cosmological framework is given as follows: inertial geodesics of the framed de~Sitter sphere are paths of massless photons; the fundamental rest-frame of massive particles is the collection of worldlines that actually describe particles which move at the rate of null-lines relative to these photons---i.e., the null-worldlines themselves; in particular, we ustilise the Hopf fibration of the 3-sphere ($SU(2)$) to identify two opposing tangent directions at each point on a 2-sphere ($SU(2)/U(1)$; with tangent vectors directed either way along the circle ($U(1)$) at each point), and define the fundamental massive particles as those which move all in the same direction at the null speed; accordingly, we write down the metric for the 3-sphere which expands in cosmic time, so that photons are now described as travelling along null-lines in that direction, as
\begin{equation}
ds^2=-B(r,t)dr^2+A(r,t)dt^2+r^2\left(d\theta^2+\sin^2\theta{d\phi^2}\right),\label{sphersymm_r_intro}
\end{equation}
where $r$ is the timelike coordinate describing the radial expansion of the actual universe in this frame; but then, by definition, the null-lines of this system must be both, the fundamentally inertial geodesics \textit{and} the inertial worldlines of particles that move through the 3-sphere at twice the fundamental null-rate in that direction, so that particles travelling along these perceived null-lines should also be identified as photons \textit{by association}; and finally, the picture is completed by assuming continuous symmetry, by which the homogeneous and isotropic universe would be filled uniformly with particles moving with all possible velocities, so that throughout the 3-sphere there will be as many particles with velocities pointing in one direction, as there are particles with velocities pointing in the other, and those velocities may in fact range in magnitude from $-\infty$ to $+\infty$.

According to a massive particle travelling with velocity between the two null-lines defined in Eq~(\ref{sphersymm_r_intro}), the paths of all other particles with similar velocities will be described as timelike worldlines. Also, photons travelling along fundamental geodesics will be incoming, while those travelling at twice the rate of a cosmic rest-frame observer will be outgoing; therefore, the coordinate $t$ in Eq.~(\ref{sphersymm_r_intro}) is associated with the proper radial direction in massive frames. This description is equivalent for particles with velocity in the opposite direction of the circle as well, which associate with the bundle of null-lines propagating in the opposite direction, thereby defining a third set of worldlines pertaining to photons (i.e., particles with twice the null-velocity in that direction). These two types of matter should therefore be quite similar, but because their relative velocities are greater than those of photons, they must also be relatively very different.

This difference is realised as follows: all of these velocities correspond to radial momenta, which enter into the statical line-elements, given by the speed-time contractions (the local form of Eq.~(\ref{SdS_pure_intro}), which is found by restricting Eq.~(\ref{Edd_contraction_intro}) to a solution of Einstein's equation), as mass-terms; the mass of a (spherically symmetric, uncharged) particle, which is identified as precisely half the horizon radius, is not a single value, but a triplet---e.g., a photon should not merely have trivial mass, which is the mass of an incoming photon in every frame, but $2m_{\gamma}=\{0,\pm\sqrt{3/\Lambda}\}$, corresponding to the three paths defined as null-lines through the two distinct massive frames; accordingly, each massive particle has either two positive and one negative mass, or two negative and one positive mass; the former are described as positive mass particles, and the latter as negative mass particles, and these are mutually invisible and gravitationally repulsive, as determined in \S~\ref{sec_DPM}.

It is significant, that each triplet, or superposition of three mass values, is as close to an objective definition of mass as can be given by the geometry: a particular value fixes the other two unambiguously; and the geometry, from which the mass is determined, pertains to this degenerate set. This can be understood kinematically as follows: due to the symmetrical properties of the description resulting from the conjunctive definitions, of actual null-lines as the fundamental rest-frame geodesics, and the actual geodesic paths of massless particles with trivial inertia as null-lines, the inertial speed of a massive particle through the universe, say $0<(d\vartheta/dt)_m<(d\vartheta/dt)_{\mathrm{null}}$ in the coordinate system of Eq.~(\ref{dS_comoving_univ_intro}), must be relatively equivalent to that of a particle with actual speed $2(d\vartheta/dt)_{\mathrm{null}}-(d\vartheta/dt)_m$. But also, from the perspective of this same cosmic rest-frame, the relative speed of a particle with `negative' mass, $-(d\vartheta/dt)_{\mathrm{null}}<-(d\vartheta/dt)_{-m'}<0$, should actually be equivalent to that of a particle \textit{moving in the opposite direction}---i.e., the same direction as `positive' mass particles---with speed $2(d\vartheta/dt)_{\mathrm{null}}+(d\vartheta/dt)_{-m'}$ greater than the null-rate prescribed for photons in that cosmic rest-frame. Then, by extension, relative speeds between $-(d\vartheta/dt)_{\mathrm{null}}$ and $-2(d\vartheta/dt)_{\mathrm{null}}$ should correspond to particles with negative mass, and those less than $-2(d\vartheta/dt)_{\mathrm{null}}$ should correspond to particles with positive mass. Accordingly, the mass of a particle is a degenerate triplet corresponding to its fundamental radial inertia (and the local SdS solutions, by utilising the relativity of inertia in the appropriate speed-time contraction, describe spacetime around particles that remain at the origin, so that this inward or outward inertia corresponds to rest-mass), and further determination than that is not possible, relatively speaking.

Therefore, the inertial velocity of a particle through the universal 3-sphere, which was used in order to determine the cosmic rest-frame for massive particles, should not be thought to correspond to a cosmological picture in which some particles `move this way', while others `move that way' through apparent space, but simply to each particle's rest-mass, as it exists in the fundamental space with non-vanishing `radial inertia', which is absorbed into the timelike coordinate in its proper speed-time contraction.

This is really the most intuitive way of describing cosmological test-particle dynamics of the fundamental Newton group, Eq.~(\ref{dS_comoving_univ_intro}); as each massive observer, with non-trivial inertial velocity through the fundamental 3-sphere (although they should all naturally consider themselves as being at rest), should naturally coordinate its metric, at any time, as   
\begin{equation}
ds^2=d\vartheta'+\sin^2\vartheta'{d\Omega}^2,
\end{equation}
so that $\vartheta'$ is the radial coordinate, which is locally like the radial coordinate of flat space, and space is conceived by such an observer, with non-trivial velocity in the (positive \textit{or} negative) $\vartheta$-direction, as infinitely many 2-spheres stacked along this dimension, containing matter which might be relatively moving through those three dimensions as well, which should each personally conceive of their own surroundings in the same way, so that each contracts their local relativistic spacetime coordinate system according to this idea of radial symmetry in proper inertial space, along with an assumption on the proper radial velocity of photons, so that the fundamental velocity of each particle must enter into the metric in another way. 

With this description in mind, we return to the cosmological picture. First of all, note that the Newton group, which is trivially described by Eq.~(\ref{dS_comoving_univ_intro}), is a 3-sphere which contracts to a finite size at $t=0$, and expands afterwards. However, by defining a `radial' timelike-coordinate in the manner specified in Eq.~(\ref{sphersymm_r_intro}), in order to contract this same group about a bundle of null-lines of de~Sitter space, the metric was constructed as one which should only describe the expanding \textit{or} contracting half of the group. Thus, the expanding cosmological description in the coordinates given by the eventual solution, Eq.~(\ref{SdS_pure_intro}), is singular at $r=0$, where expansion begins, whereas the contracting half collapses to this singularity. Interestingly enough, the mass in the cosmological solution actually is a single, well-defined~value.

Now, in the proposed large-scale scenario, this expanding universe is supposed to begin with a homogeneous distribution of positive and negative mass particles, which naturally dissociate, through mutual gravitational repulsion, into a sort of lattice distribution of clusters of matter with positive and negative mass, which are mutually invisible. This repulsion will not, however, contribute to the cosmic rate of expansion, which is a well-defined geometric property. Instead, the gravitational repulsion that would be dynamically felt within any massive cluster, due to the presence of an effective local shell of `negative' mass, would be equivalent to an additional centrally attractive force within the cluster.

This theory for the true nature of locally non-interacting dark matter and natural formation of the void-filament structure observed in our Universe, as well as observed baryon asymmetry, is quite similar to a recent proposal by Massimo Villata \cite{Villata2011a,Villata2011b}, except that he has suggested that the repulsive gravitational interaction between matter and anti-matter should contribute to the cosmic expansion rate. In fact, the theory is even consistent with Villata's idea that these anti-matter particles should, in one way, be describable as moving backwards in time, even though we now also describe everything as remaining within the three-dimensional universe, moving forward in its absolute cosmic time.

And as this occurs, the global expansion rate of this RW universe should be determined from the cosmological SdS solution, Eq.~(\ref{SdS_pure_intro}). For the universe really is, by construction, homogeneous and isotropic; and although it is not synchronous in the proper coordinate frame of massive particles in the cosmic rest-frame, it must appear to them as such, because those non-trivial fundamental worldlines do fan out isotropically as the universe evolves, so that relative to any one, all others would appear to be isotropically receding. Furthermore, the apparent rate of cosmic expansion must be given according to the change in $r$, as measured in the proper time coordinate system of the fundamental rest-frame, Eq.~(\ref{SdS_scalefac_intro}), which is the same as the flat $\Lambda$CDM scale-factor.

The connection to FLRW theory is then trivial: in that theory we assume a universe which is isotropic and homogeneous, and expands in absolute time at a rate which is determined by three parameters---the density and pressure of the various forms of matter it contains, and its spatial curvature; the form of the scale-factor which describes this expansion is constrained, through Einstein's equation, after assuming perfect fluidity of matter as an acceptable large-scale approximation, to a solution of the two Friedman equations, which depend on these three physical parameters; we then solve this under-determined system of equations by considering equations of state, relating pressure and density for various types of perfect fluid that the universe might contain, and subsequently model observables in order to constrain the form of the scale-factor, which depends on the values of the three physical parameters; accordingly, observations have constrained the scale-factor quite well to the form of Eq.~(\ref{SdS_scalefac_intro}); the theoretical universe described in this thesis must bear that very same appearance, since it appears from the fundamental rest-frame to be expanding isotropically, and scaling, according to the proper measure of the expansion factor in that frame, according to Eq.~(\ref{SdS_scalefac_intro}).

Some further points to consider are: the expansion occurs for purely geometrical reasons, rather than due to the local dynamics of material in causally connected regions of the universe, so there is no horizon problem; the apparent scale-factor is precisely that of the flat $\Lambda$CDM universe for all time, so there is no flatness problem because any measurement of cosmic expansion must constrain the curvature parameter to zero; the prior resolution of both of these problems, together with the fact that the rate of evolution is well-defined for all time, means that there is no need for an inflationary epoch; in fact, because the cosmic expansion is not driven by the local energy content, but is a well-defined prior property of the universe and the underlying void, there is also no cosmological constant problem, because the vacuum expectation value from quantum theory simply cannot contribute to the expansion rate.

But beyond that, the geometry of this universe is essentially similar to the one that describes a locally spherically symmetric mass distribution. Furthermore, as shown in the thesis, the radial coordinate in the local geometry actually should not exist on the supposed `timelike intervals', so that causally coherent, spherically symmetric gravitational collapse beyond the horizon radius, which cannot occur until the end of cosmic time, must not be described by that solution. Instead, by hypothesising that spherically symmetric gravitational collapse in one universe should, at the end of its cosmic time, produce a contracting Newton group, the picture is completed as one which is not merely consistent with cosmology, but should actually provide for itself the basic cause for the cosmic time that needs to be assumed, because this new universe might also be described by the contracting cosmological SdS solution, with particle dynamics described in the same way, which would eventually reach $r=0$ in those coordinates, at the equator of the de~Sitter sphere that was framed at the final stage of local collapse in the progenitor universe, and then emerge from that singularity, in a hot state following complete gravitational collapse, to be described now by the expanding cosmological SdS solution. Thus, a connection is made between big bang universes like ours and gravitational collapse, which should finally account for becoming in this theory.

And so, the thesis, which is sometimes deeply speculative, should not be accepted purely according to its philosophical appeal---viz., as it is very intuitive and common sensical, and admits only the most reasonable and globally consistent possibilities when confronted with dilemmas,---but because this leads to a metrically well-defined, fundamentally covariant theory of Physical Reality which does not just fix, but essentially resolves the most significant problems of the current dynamical theory, providing logically consistent, realistic answers for the cause of the Big Bang and the basic Nature of Time through causally coherent gravitational collapse, according to a theory which does not require any censorship or protection conjectures to resolve paradoxes of black hole creation in finite cosmic time or time travel to a past that does not exist, and leads naturally to a theory for the origin and evolution of the large-scale structure we observe in the Universe, in a continuously symmetric cosmology describing a universe that does indeed expand for all time purely due to $\Lambda$, with evolution which would be observed to have taken place by all fundamental observers precisely according to the flat $\Lambda$CDM model that has been empirically constrained.\\\\\\








\noindent{The development of this theory is presented as follows. The first chapter is an examination of the historical development and current status of modern cosmology, beginning from the initial evidence for cosmic expansion and the theoretical modelling that became possible with the advent of general relativity theory; i.e., the era of modern cosmology is defined as having begun in the years 1912 -- 1917. Special attention is paid to the theories surrounding $\Lambda$, with an eye kept to the development of a new theory which would be consistent with the principles used to establish RW spacetime as a cosmological background geometry, but also consistent with Eddington's interpretation, that cosmic expansion might be objectively required by a geometrical property of Physical Reality. Therefore, after outlining the details in the history of early modern cosmology, from which the RW line-element emerged as providing the general description of an isotropic and homogeneous expanding universe (\S~\ref{sec_EMC}), the reasoning that was given for the two competing interpretations of expansion is discussed in \S~\ref{sec_ICE}, in order to establish the case for Eddington's interpretation as a more realistic alternative to Einstein's. Then, various ways in which the cosmological constant was subsequently incorporated are discussed in \S~\ref{sec_TSC}, as these eventually led to the interpretation of a dark energy component of the Universe after its existence was finally discovered in 1998 \cite{Riess1998,Perlmutter1999}.}

Now, the big bang FLRW models are not only inconsistent with the interpretation of $\Lambda$ as the basic reason for cosmic expansion, as they do lend themselves more naturally to the description of $\Lambda$ as a dark energy component of a universe that resulted from a singularity (unless the theory of gravity itself is modified in order to avoid this), but there are other significant problems, such as the observed flatness, isotropy on the largest scales---which, together with the cosmological principle, implies large-scale homogeneity (isotropy from every point), but also a horizon problem since this could not result from any causal influence in the universe that is described just so,---and the fact that all the matter we know anything about, which we can directly observe, makes up only four percent of the measured total energy in the Universe. Actually, the most significant of the cosmological problems is likely the paradoxical implication from relativity theory, that the Universe may not be evolving at all, but that each present instant might be instead just one subjective slice of a four-dimensional Block Universe.

But the recurring theme here, especially in the more well-known problems in cosmology (horizon, flatness, dark matter/dark energy), is not that there is any discrepancy between the observational data and the physical model, but that the empirically constrained parameters are inconsistent with our expectations of what should realistically be. Therefore, the remainder of the thesis is an attempt to revisit those expectations, and derive a theory from rationalised first principles which would be consistent with the evidence: that the Universe is expanding, and probably due to a cosmological constant, with a measurable expansion rate equivalent to the description in which the (other) dominant energy component indicates the presence of an appreciable amount of invisible dust; that it began expanding from a hot initial state; and that it must appear perfectly isotropic and flat to all observers. Thus, the theory would be observationally equivalent to the standard (flat $\Lambda$CDM) concordance model, but would provide reasons why things ought to be the way they appear.

This begins, in Chapters~\ref{chap_AER} and \ref{chap_DIRT}, with an examination of the physical nature of time, which is essential to the problem of deriving a nontrivial relativistic cosmological model. Therefore, these chapters mainly present an argument for why relativity theory in general can and must be interpreted as describing the evolution of a three-dimensional universe, in a manner that is consistent with, but generalises the trivial prescription of an absolute cosmic time that is made in deriving the RW line-element; thus, e.g., it is meaningful to say in general, that a particle which has existed in the Universe since the Big Bang has \textit{really} `been around' for the past 13.7 billion years, even though it may have been travelling all that time through the cosmic rest-frame at some significant fraction of the speed of light so that its relative proper time would be considerably less; i.e., the purpose of this argument is to generally reconcile relativity theory with the sense of the real present that relativistic cosmology takes for granted.

So, beginning with an investigation into the problem of the physical nature of time in its historical context, the argument proceeds from an epistemic standpoint. The purpose of Chapter~\ref{chap_AER} is therefore to discuss the foundational insights that were made in the philosophy of time, from Heraclitus and Parmenides, through the Atomists, Aristotle, and St.~Augustine, eventually to Newton and Einstein, and to discriminate between the epistemological methods that developed along the way, in order to form a case against the unnatural inferences that were finally made by Einstein and Minkowski. As such, when discussing the theories of Heraclitus (\S~\ref{sec_OH}) and Parmenides (\S~\ref{sec_OE}), along with the reasoning behind them, a great deal of consideration is given to the epistemological methods that they used; and this eventually leads, in \S~\ref{sec_ANE}, through further discussion of how those methods were incorporated into those of their successors, to the formulation of the natural epistemic stance that is taken in this thesis. In \S~\ref{sec_OAA}, the Atomists' theory is presented as one based on a reduction of Parmenides' argument, and then Augustine's recognition that time must be equable duration (as Newton eventually put it), is discussed in \S~\ref{sec_OT}, before finally coming to the interpretations of special relativity theory by Einstein and Minkowski, in \S~\ref{sec_OIS}, and their abandonment of the significant philosophical inferences that had previously been made with regards to the Nature of Time.

The chapter therefore sets up the argument, that Einstein, by simply rejecting Newton's absolute time, space, and motion, and taking the positivist stance regarding the relativity of simultaneity, but the stance of a realist when interpreting the consequential meaning of that inference, was compelled to agree with Minkowski that we must live in an absolute world like that of Parmenides, when he should rather have reduced Newton's philosophically well-motivated theory to one that would be consistent with the causal and intertial structures of the spacetime metric in special relativity theory.

Chapter~\ref{chap_DIRT} then proceeds with an attempt to reconstruct relativity theory as a truly dynamical one, consistent with the principles of Heraclitus, the Atomists, Augustine, Galileo, and Einstein, with the idea that there should be a four-dimensional void, acting as a general relativistic \textit{F\"{u}hrungsfeld} through which the three-dimensional universe should evolve, with events in spacetime emerging with the necessary metrical structure. Therefore, the past and future events of spacetime are described as being \textit{ideal} occurrences in a truly evolving three-dimensional present. The case for special relativity is neatly illustrated in \S~\ref{sec_OJGKD}, with the example of a barograph in which a number of pens exist in a pre-arranged line with their paths traced out on graph paper that moves underneath by clockwork; and a reassessment of the objectively real potentialities that might give rise to the same empirical model (viz.\ special relativistic spacetime) is presented in \S~\ref{sec_ARP}, with the eventual conclusion that this alternative scenario should have the best chance of leading to a determination of the essential nature of time. At the end of this section, an argument against the causally coherent formation, and present existence of black holes in our Universe begins to develop. Then, in \S~\ref{sec_RPT}, the curvature and geometrodynamics of spacetime are considered, and it is eventually hypothesised that these might actually be described as relatively projective properties, purely due to particle dynamics in de~Sitter space, rather than as the result of a dynamically warping \textit{F\"{u}hrungsfeld}. 

The new cosmological theory really begins to build from this point; and, in the latter half of \S~\ref{sec_RPT}, upon the assumption that the absolute cosmic time should be given according to the induced metric of a de~Sitter sphere, as an absolute background structure with intrinsic Lorentzian signature, two line-elements---one local, and one cosmological---emerge from considerations of how the events of four-dimensional spacetime might be algebraically coordinated. These are the SdS metrics---and the two physically different, yet algebraically equivalent statical geometries are described in \S~\ref{sec_GPSMRF}, using dimensionless coordinates normalised by the cosmical quadric of curvature, according to their fundamental parametrisation. In \S~\ref{sec_DPM}, the description of physical mass in the geometry is considered, and a possible implication for the nature of dark matter and the formation of large-scale structure in the Universe is discussed. Then, in \S~\ref{sec_MC}, the cosmographical metric is analysed, and it is shown that the factor of expansion in the essentially isotropic and homogeneous universe would be perceived, by any observer in the fundamental rest-frame defined by a bundle of worldlines satisfying Weyl's principle, to evolve in proper time precisely according to the flat $\Lambda$CDM scale factor. 

It is significant to note---with reference to my basic hypothesis, that a different cosmological ontology should be able to be given so that we might rationally \textit{expect} the observation of flat $\Lambda$CDM-type expansion,---that in this scenario the bundle of fundamental worldlines are not orthogonal to the isotropic and homogeneous cosmic hypersurfaces, which is one of the basic assumptions that Robertson did make, and is therefore one of the basic mathematical principles of FLRW cosmology \cite{Rugh2010}. In fact, the fundamental worldlines in this theory are not geodesics of the fundamental background structure that is assumed---and they are not even timelike, but null;---rather, the bundle of timelike geodesic worldlines that are orthogonal to the cosmic hypersurfaces (in which frame the comoving spacetime line-element \textit{is} RW) should actually be massless, e.g.~with $M=0$ in their corresponding local statical (S)dS geometries.

Finally, because this cosmological model, which begins from a coordinate singularity of the relevant metric, is described only on the expanding half of the de~Sitter sphere, it is pointed out, in \S~\ref{sec_CBHO}, that because the corresponding description on the contracting half of the sphere is relevant to the description of the final stage of spherically symmetric gravitational collapse, it may be possible to describe the emergence of such a universe as the result of gravitational collapse in a prior one, which process would give rise to the absolute time that has to be assumed in the model.